\clearpage
\section{The explicit dyadic arithmetic tower}\label{app:dyadic_tower}

No character, quotient-dimension,
source attachment, pair-factor, or universal Chapter XI closure is
asserted here.

\hypertarget{dyadic_tower-statement}{%
\subsection{1. Statement}\label{dyadic_tower-statement}}

Let F=F\_2((t)), let 0\textless h\_1\textless\ldots\textless h\_d be odd
positive integers, and set

\begin{Verbatim}[fontsize=\small,breaklines=true,breakanywhere=true]
q_i=3h_i,  N=2^d.
\end{Verbatim}

The explicit recursion below constructs B\_d in M\_N(F) such that

\begin{Verbatim}[fontsize=\small,breaklines=true,breakanywhere=true]
F[B_d] is a totally ramified maximal field normalizing
   the complete period-N lattice chain,
v_A(B_d)=-2^(d-1)q_d,
k_0(B_d,A)=-q_1,
tr_F(B_d)=1.
\end{Verbatim}

There are explicit native constructors of ranks 2\^{}j for the largest j
exponents. Their identity-amplifications give the actual defining
sequence with field degrees 2,4,\ldots,2\^{}d and critical levels

\begin{Verbatim}[fontsize=\small,breaklines=true,breakanywhere=true]
n_j=2^(d-j)q_(d-j+1),  1<=j<=d.
\end{Verbatim}

Every parameter in this ambient defining sequence has the same leading
level n\_1. The quantities n\_2,\ldots,n\_d are its successive critical
approximation levels, not new leading levels of the ambient matrices.

Every native constructor has trace 1. Every correction from an
identity-2 amplification of the preceding native constructor to the next
native constructor also has trace 1. In full ambient rank 2\^{}d, the
earlier amplified approximations instead have trace 0, because their
multiplicities are even. This distinction is essential.

\hypertarget{dyadic_tower-the-minimal-degree-two-building-block}{%
\subsection{2. The minimal degree-two building
block}\label{dyadic_tower-the-minimal-degree-two-building-block}}

Let L be a complete discretely valued field of characteristic 2 with
residue field F\_2, and choose a uniformizer t\_L. For odd
q\textgreater0 put a=(q+1)/2 and define

\begin{Verbatim}[fontsize=\small,breaklines=true,breakanywhere=true]
E=L(alpha),  alpha^2-alpha=t_L^(-q),
z=t_L^a alpha.
\end{Verbatim}

The Artin--Schreier polynomial is irreducible: an element x\^{}2-x of
negative L-valuation has even negative valuation. Its derivative is 1,
and the valuation equation forces ramification index 2. Thus E/L is
separable and totally ramified, v\_E(alpha)=-q, and alpha is minimal.

The element z has E-valuation 2a-q=1 and satisfies

\begin{Verbatim}[fontsize=\small,breaklines=true,breakanywhere=true]
f(z)=z^2+t_L^a z+t_L=0.
\end{Verbatim}

Thus 1,z is a uniformizer power basis, and

\begin{Verbatim}[fontsize=\small,breaklines=true,breakanywhere=true]
alpha=z/t_L^a,  Tr_(E/L)(alpha)=1.
\end{Verbatim}

\hypertarget{dyadic_tower-an-explicit-normalized-corestriction-and-lift}{%
\subsection{3. An explicit normalized corestriction and
lift}\label{dyadic_tower-an-explicit-normalized-corestriction-and-lift}}

On E use the coefficient basis 1,z. Its trace-dual elements are

\begin{Verbatim}[fontsize=\small,breaklines=true,breakanywhere=true]
z_0^*=(z+t_L^a)/t_L^a,
z_1^*=1/t_L^a.
\end{Verbatim}

Let pr(T\_ij)=z\^{}i z\_j\^{}* be the trace corestriction. Put

\begin{Verbatim}[fontsize=\small,breaklines=true,breakanywhere=true]
kappa=t_L^a/z=alpha^(-1),
s=kappa pr,
u_0=s(P0)=t_L/z^2.
\end{Verbatim}

Here P0 is constant-coefficient projection, applied separately to each
coordinate when acting on E\^{}m. The identity for u\_0 follows from
f(z)=0. It is a unit: v\_E(t\_L)=v\_E(z\^{}2)=2. More explicitly,

\begin{Verbatim}[fontsize=\small,breaklines=true,breakanywhere=true]
s(T_i0)=z^i(1+t_L^a/z),
s(T_i1)=z^(i-1).
\end{Verbatim}

Since v\_E(t\_L\^{}a/z)=2a-1=q\textgreater0, these preserve the expected
leading grades. On End\_L(E\^{}m), take the complete interleaved chain
in the basis w\_(mi+b)=z\^{}i v\_b. Its centralizer is M\_m(E) with its
complete E-chain. The formulas give s(P\^{}r) contained in Q\^{}r for
all integers r.

For C in M\_m(E), define the explicit L-linear operator

\begin{Verbatim}[fontsize=\small,breaklines=true,breakanywhere=true]
Lift_E(C)=u_0^(-1) C P0=(z^2/t_L) C P0.
\end{Verbatim}

B-bimodularity gives s(Lift\_E(C))=C. The lift preserves exact
valuation: u\_0\^{}(-1) is an E-unit, and expansion in 1,z gives
coefficient valuations distinct modulo 2. Each coefficient entry has the
same grade in the interleaved L-chain and the E-chain. In particular
\(s(P^r)=Q^r\).

\hypertarget{dyadic_tower-exact-trace-transport}{%
\subsection{4. Exact trace transport}\label{dyadic_tower-exact-trace-transport}}

For every C in M\_m(E), ordinary matrix traces satisfy

\begin{Verbatim}[fontsize=\small,breaklines=true,breakanywhere=true]
tr_L(Lift_E(C))
   =Tr_(E/L)(z_0^* u_0^(-1) tr_E(C))
   =Tr_(E/L)(kappa^(-1) tr_E(C))
   =Tr_(E/L)(alpha tr_E(C)).                 (T2)
\end{Verbatim}

The first equality follows by taking the L-trace of each rank-one
operator yP0, which is the constant coefficient of y. This coefficient
is Tr\_(E/L)(z\_0\^{}* y).

Thus the trace of a lift is 1 whenever the inner ordinary E-matrix trace
is 1. It is not asserted that arbitrary trace values are unchanged.

\hypertarget{dyadic_tower-recursive-actual-matrices-and-fields}{%
\subsection{5. Recursive actual matrices and
fields}\label{dyadic_tower-recursive-actual-matrices-and-fields}}

Define B(L,t\_L;h\_1,\ldots,h\_d) as follows. For d=1 use the
multiplication matrix of alpha from Section~2 with q=q\_1. For
d\textgreater=2 apply Section~2 with q=q\_d, set m=2\^{}(d-1),
recursively construct

\begin{Verbatim}[fontsize=\small,breaklines=true,breakanywhere=true]
C=B(E,z;h_1,...,h_(d-1)),
\end{Verbatim}

and put

\begin{Verbatim}[fontsize=\small,breaklines=true,breakanywhere=true]
B(L,t_L;h_1,...,h_d)=alpha I_m+(z^2/t_L) C P0.       (D2)
\end{Verbatim}

This is a finite algebraic recipe. It uses specified separable quadratic
extensions and no unknown support, genericity choice, or correspondence
inverse.

The degree and simplicity proof is by induction. The base alpha is
minimal and generates E. At a later step let gamma=alpha I\_m and
b=Lift\_E(C). The outer minimal stratum has

\begin{Verbatim}[fontsize=\small,breaklines=true,breakanywhere=true]
n=2^(d-1)q_d,  k_0(gamma,A)=-n.
\end{Verbatim}

The inner induction hypothesis gives an actual maximal field E{[}C{]}
normalizing the complete E-chain with

\begin{Verbatim}[fontsize=\small,breaklines=true,breakanywhere=true]
R=2^(d-2)q_(d-1),  v(C)=-R,  k_0(C)=-q_1.
\end{Verbatim}

Because h\_d\textgreater h\_(d-1), we have n\textgreater2R. Therefore

\begin{Verbatim}[fontsize=\small,breaklines=true,breakanywhere=true]
-2R > -n-q_1 = k_0(gamma,A)+k_0(C).
\end{Verbatim}

The derived stratum with parameter s(b)=C is pure, and its field is
maximal. Bushnell--Kutzko Remark 2.2.7 consequently gives the field
generated by the actual gamma+b, not merely an equivalent replacement.
It has ramification and degree 2\^{}d, residue degree one, normalizes
the complete chain, and has k\_0=max(-n,-q\_1)=-q\_1. Since
v(b)\textgreater v(gamma), its level is n and {[}A,n,0,gamma+b{]} is
simple.

The derived model E{[}C{]} is not identified with L{[}gamma+b{]}. In
particular no literal embedding of E into the latter field is claimed.

\hypertarget{dyadic_tower-trace-1-at-every-native-stage}{%
\subsection{6. Trace 1 at every native
stage}\label{dyadic_tower-trace-1-at-every-native-stage}}

The base matrix has trace Tr\_(E/L)(alpha)=1. Inductively tr\_E(C)=1, so
(T2) gives tr\_L(Lift\_E(C))=1. For d\textgreater=2 the L-trace of alpha
I\_m is

\begin{Verbatim}[fontsize=\small,breaklines=true,breakanywhere=true]
m Tr_(E/L)(alpha)=0,
\end{Verbatim}

because m is even. Hence tr\_L(B)=1 at every native rank.

For any j\textgreater=1, let A\_j denote the native rank-2\^{}j
constructor using a fixed list of the largest j exponents. The next
native rank-2\^{}(j+1) constructor A\_(j+1) uses the same largest
exponents and adds the next smaller one. With their compatible
interleaved bases, put

\begin{Verbatim}[fontsize=\small,breaklines=true,breakanywhere=true]
Delta_(j+1)=A_(j+1)-A_j tensor I_2.
\end{Verbatim}

Then

\begin{Verbatim}[fontsize=\small,breaklines=true,breakanywhere=true]
tr(Delta_(j+1))=1-2*1=1                    in characteristic 2.
\end{Verbatim}

This native trace identity is stronger and more relevant than a claim
obtained merely by amplifying to a rank divisible by two.

\hypertarget{dyadic_tower-full-arithmetic-defining-sequence}{%
\subsection{7. Full arithmetic defining
sequence}\label{dyadic_tower-full-arithmetic-defining-sequence}}

For the final rank 2\^{}d define

\begin{Verbatim}[fontsize=\small,breaklines=true,breakanywhere=true]
B^[j]=B(F,t;h_(d-j+1),...,h_d) tensor I_(2^(d-j)).
\end{Verbatim}

All truncations use the same top q\_d to select the first extension and
uniformizer, the same next q\_(d-1) to select the next extension, and so
on. The coefficient projection and lift commute with amplification. Thus
these are literal truncations of (D2).

Their parameter fields have degree 2\^{}j, and their critical levels in
the ambient complete chain are n\_j=2\^{}(d-j)q\_(d-j+1).
Valuation-preserving lifts give

\begin{Verbatim}[fontsize=\small,breaklines=true,breakanywhere=true]
v_A(B^[j+1]-B^[j])=-n_(j+1).
\end{Verbatim}

Later differences have strictly greater valuation, since
n\_j\textgreater2n\_(j+1). This proves that the displayed parameters
form an actual defining sequence of simple approximations.

Their ambient traces are

\begin{Verbatim}[fontsize=\small,breaklines=true,breakanywhere=true]
tr_F(B^[j])=0 for j<d,  tr_F(B^[d])=1.
\end{Verbatim}

Consequently ambient corrections before the last step have trace 0, and
the last correction has trace 1. Every one of these trace values is
integral. On a scalar principal unit u, the additive correction
character has value psi\_F((u-1)tr(delta)). This is 1 for an additive
character trivial on the maximal ideal. No assertion about full simple
character extension follows here without its separate intersection and
transfer argument.

\hypertarget{dyadic_tower-concrete-rank-eight-instance-and-limits}{%
\subsection{8. Concrete rank-eight instance and
limits}\label{dyadic_tower-concrete-rank-eight-instance-and-limits}}

For h=(1,3,5), the construction gives actual approximation-field degrees
2,4,8 and critical levels 60,18,3. The supplementary exact program
constructs its Laurent polynomial matrices, verifies the normalized lift
identities, and checks the native/ambient trace distinctions and matrix
grades. It does not prove field existence or simplicity by bounded
computation: those are consequences of Section~5.

The odd-prime Heisenberg quotient and conductor counting from Block166
cannot simply be copied: q\_1 is odd here. This note makes no dyadic
quotient, full-character, stabilizer, phase, source-attachment, or H5
claim. It also changes lower parameter conjugacy classes, as in the
corresponding odd-prime Artin--Schreier construction.

\hypertarget{dyadic_tower-dependencies-and-provenance}{%
\subsection{Dependencies and
provenance}\label{dyadic_tower-dependencies-and-provenance}}

Primary theorem: Bushnell--Kutzko, The admissible dual of GL(N) via
compact open subgroups, Remark 2.2.7 (printed p. 55), together with the
amplification normalization and defining-sequence conventions already
checked in Block166.

Retained primary source text SHA256:
\path{d69db0827256707a34cd2562af68b9de11fc889ab3f43535ec47b620d58ca9c8}.

Block166 inputs are pinned separately in \path{INPUT_PINS167.json} and are not
modified by this continuation. The Artin--Schreier, trace-dual, and
trace calculations needed specifically for p=2 are proved above.
